\section{Uniform original normal ranks on the genuine common open}

Let $\mathfrak I$ be an original coherent whole ideal sheaf in the split
ambient coordinates, and let $I_a$ be its actual translated original whole
germ ideal at the point $(a,0)$. The common coefficient construction gives
a finite type domain $D$, a finite family of true homogeneous relations,
and their whole generated ideal $H\subset D[Y_1,\ldots,Y_t]$, together
with genuine coefficient homomorphisms $\rho_a$ whose whole ideal images
are exactly the original initial ideals of $I_a$. Let $b\in D\setminus\{0\}$
be the actual uniform leading denominator of this one whole ideal, and
write $L=D[b^{-1}]$. These choices are independent of all degrees and
truncation orders. The source-facing objective is the whole local-ring
Hilbert--Samuel quotient in Bierstone--Milman's ``Use of the Hilbert--Samuel
function''. The following argument supplies its relative-filtration rank and
literal residue-tensor dimensions first; the geometric maximal-ideal quotient
comparison and descent remain separate producers.

\begin{lemma}[Actual coefficient-image standard monomial bases]
\label{hir:relative-normal-nearby-coefficient-standard-basis}
\source{bierstone-milman08-functoriality}
\lean{Hironaka.HironakaRelativeNormalNearbyUniformRanks.coefficientImageStandardMonomialBasis,Hironaka.HironakaRelativeNormalNearbyUniformRanks.coefficientImageStandardMonomialBasis_apply}
\leanok
\uses{hir:relative-normal-uniform-localized-whole-basis,hir:relative-normal-coefficient-polynomial-quotient-base-change}
Fix the actual field-standard exponent family $E$ of $H$, computed over
the fraction field of $D$. For every coefficient homomorphism $f:D\to T$
with $f(b)$ a unit, the whole polynomial quotient by the image of $H$ has
a $T$-basis indexed by this same $E$. Its vector indexed by $\alpha$ is
the actual quotient class of $Y^\alpha$ with coefficient $1$.
\end{lemma}
\begin{proof}
The unit condition gives a genuine homomorphism $L\to T$. The existing
uniform standard monomial basis of $L[Y]/HL[Y]$ has its literal monomial
numerator formula. Tensor that basis with $T$ over $L$. The whole polynomial
coefficient base-change algebra isomorphism sends $1\otimes[Y^\alpha]$
to the actual class of $Y^\alpha$ in $T[Y]$ modulo the extended whole
ideal. The composition of $D\to L\to T$ is $f$, so this extended ideal
is exactly the whole image of $H$ under $f$. Tensoring a basis and applying
these genuine isomorphisms proves both the asserted basis and its actual
numerator formula, without requiring $f$ to be injective.
\end{proof}

\begin{lemma}[Actual homogeneous degree fibres of a whole monomial basis]
\label{hir:relative-normal-nearby-standard-degree-basis}
\source{bierstone-milman08-functoriality}
\lean{Hironaka.HironakaRelativeNormalNearbyUniformRanks.originalDegreeCoefficientImageMonomialBasis,Hironaka.HironakaRelativeNormalNearbyUniformRanks.originalDegreeCoefficientImageMonomialBasis_finrank}
\leanok
\uses{hir:relative-normal-nearby-coefficient-standard-basis,hir:relative-normal-intrinsic-associated-quotient-algebra-dictionary}
For any actual nearby original quotient dictionary, the preceding whole
basis identifies each literal original degree-$q$ layer with the finite
free coefficient module indexed by
$E_q=\{\alpha\in E:|\alpha|=q\}$. Thus its rank is $\#E_q$.
\end{lemma}
\begin{proof}
The intrinsic original quotient algebra dictionary sends each whole basis
vector to the actual degree-$|\alpha|$ original monomial symbol. For a
coefficient function supported in $E_q$, include it in the whole exponent
family, apply the whole basis isomorphism, and project to degree $q$.
Its original degree inclusion recovers the whole vector, because each
basis vector used in the sum already lies in degree $q$. This gives
injectivity. Conversely, include any literal degree-$q$ vector in the
whole direct sum and take its unique whole basis coefficients. Projection
to degree $q$ kills every basis vector of other degree and retains every
basis vector of degree $q$. Uniqueness of coefficients therefore forces
the original coefficient function to be supported in $E_q$. This gives
surjectivity and the claimed genuine degree isomorphism. Every exponent
of degree $q$ has all its entries at most $q$, so $E_q$ injects into a
finite set of bounded normal exponents. Its cardinal is therefore finite,
and the degree isomorphism gives the asserted rank.
\end{proof}

\begin{theorem}[One original-input open and one complete original rank sequence]
\label{hir:relative-normal-nearby-all-degree-uniform-rank}
\source{bierstone-milman08-functoriality}
\lean{Hironaka.HironakaRelativeNormalNearbyUniformRanks.exists_original_nearby_all_degree_uniform_rank_open}
\leanok
\uses{hir:relative-normal-original-common-domain-coefficient-initial-family,hir:relative-normal-nearby-standard-degree-basis}
Original coherence alone supplies one open set $U$ in the actual base,
meeting every origin neighbourhood, and one sequence $r(q)\in\mathbb N$,
such that for every $a\in U$ and every $q$, the actual original layer
$L_q(I_a)$ is finite free over the actual base germ ring and has rank
$r(q)$. The same $U$ works for all $q$; it need not contain the origin.
\end{theorem}
\begin{proof}
Apply the original-input common-domain relation producer to obtain the
actual $D,H,\rho_a$ above on its common neighbourhood $W$. Represent the
one nonzero denominator $b$ by its true coefficient function $g$.
Its original germ is nonzero, because $D$ was the quotient by the complete
origin kernel. Set $U=W\cap\{g\ne0\}$. Holomorphicity makes $U$ open.
If some origin neighbourhood failed to meet $U$, $g$ would vanish on
its intersection with $W$, contradicting that nonzero original germ.
For $a\in U$, the actual translated germ $\rho_a(b)$ evaluates to
$g(a)\ne0$, hence is a unit. Apply the actual coefficient-image monomial
basis and the whole ideal equality supplied internally by the common-domain
producer. Apply the preceding actual homogeneous degree dictionary to the
original translated ideal $I_a$. Its degree-$q$ index family is always
the same finite set $E_q$, independently of $a$. Taking $r(q)=\#E_q$
proves all assertions with this single $U$.
\end{proof}

\begin{theorem}[Every actual whole original truncation has the common rank sum]
\label{hir:relative-normal-nearby-all-truncation-uniform-rank}
\source{bierstone-milman08-functoriality}
\lean{Hironaka.HironakaRelativeNormalNearbyUniformRanks.exists_original_nearby_all_truncation_uniform_rank_open}
\leanok
\uses{hir:relative-normal-nearby-all-degree-uniform-rank,hir:relative-normal-uniformtrunc-original-short-exact}
For the same original-input open and the same sequence $r$, every actual
whole original power truncation $R/(I_a+J^s)$ is finite free over the
base germ ring and has rank $\sum_{q<s}r(q)$, for all $a\in U$ and all
$s\in\mathbb N$.
\end{theorem}
\begin{proof}
At $s=0$ the ideal contains $J^0=R$, so the quotient is zero and its
rank is the empty sum. At the successor step, the literal original
power filtration gives the actual short exact sequence with left term
$L_s(I_a)$, middle term $R/(I_a+J^{s+1})$, and right term
$R/(I_a+J^s)$. The right term is finite free by induction and therefore
projective, so this genuine sequence splits. The middle module is
therefore the direct sum of the two finite free endpoint modules.
Their ranks add, giving $r(s)+\sum_{q<s}r(q)=\sum_{q<s+1}r(q)$.
No new neighbourhood or denominator is chosen in this induction.
\end{proof}

\begin{theorem}[Complete dimensions of the actual original algebraic fibre powers]
\label{hir:relative-normal-nearby-complete-fibre-power-dimensions}
\source{bierstone-milman08-functoriality}
\lean{Hironaka.HironakaRelativeNormalNearbyUniformRanks.exists_original_nearby_all_fibre_truncation_dimension_open}
\leanok
\uses{hir:relative-normal-nearby-all-truncation-uniform-rank}
Specialize the actual base germ ring at its origin residue homomorphism to
$\mathbb C$. On the same $U$, the complex dimension of every literal
fibre module $\mathbb C\otimes_B R/(I_a+J^s)$ is
$\sum_{q<s}r(q)$. These are dimensions of the genuine algebraic fibres
of the whole original normal-power quotients.
\end{theorem}
\begin{proof}
Each actual truncation has a finite basis with the preceding common
cardinal. Tensor this true basis with $\mathbb C$ using the actual base
evaluation homomorphism. The resulting basis of the genuine tensor fibre
has the same index family, so its complex dimension equals that same
cardinal. The already proved common truncation rank therefore gives the
asserted full sequence. Identification with the geometric fibre whole
maximal-ideal Hilbert--Samuel quotient requires the separate literal
base-maximal-ideal extension and ideal-power comparison; no geometric
Hilbert--Samuel descent is asserted by this tensor dimension statement.
\end{proof}
